\section{Regularization and multiplication of fibre bounds}\label{reg:section}

We first record two finite-tree facts.  Neither requires lower density
estimates for the original measure.  Dyadic cubes are taken half open, so
that they form genuine measurable partitions; a fixed translation and
rescaling places each measure under consideration in one unit cube.
Related decompositions into dyadically regular measures occur in
\cite[Corollary~3.5]{KS2019}; we prove the version needed here.

\begin{lemma}[Regular classes]\label{reg:regularization}
Let $n\ge1$, let $\rho$ be a probability on a unit dyadic cube in
$\R^n$, and put $R=2^N$.  For every fixed $0<\lambda<1$ and all
sufficiently large $N$, there are pairwise disjoint sets $E_\alpha$,
each a union of depth-$N$ cubes, with the following properties:
\begin{enumerate}
\item Their number is at most $N^{C_n/\lambda}$, each has
      $\rho(E_\alpha)\ge R^{-\lambda}$, and
      \[
      \rho\Bigl(\R^n\setminus\bigcup_\alpha E_\alpha\Bigr)
      \le R^{-2}+N^{C_n/\lambda}R^{-\lambda}
      =R^{-\lambda+o(1)}.
      \]
\item For $\rho_\alpha=\rho|_{E_\alpha}/\rho(E_\alpha)$ there is a
      nondecreasing $n$-Lipschitz function $m_\alpha:[0,1]\to[0,n]$,
      with $m_\alpha(0)=0$, such that every occupied depth-$j$ cube
      satisfies
      \begin{equation}\label{reg:cell-mass}
      R^{-m_\alpha(j/N)-C_n\lambda}
      \le\rho_\alpha(Q)\le
      R^{-m_\alpha(j/N)+C_n\lambda},\qquad 0\le j\le N.
      \end{equation}
\item If $\rho(B(x,r))\le Kr^s$ for $R^{-1}\le r\le1$, with
      $0\le s\le n$ and $K\ge1$, then
      \begin{equation}\label{reg:profile-lower-bound}
      m_\alpha(u)\ge su-C_n\lambda-\log_R K,
      \qquad 0\le u\le1.
      \end{equation}
\end{enumerate}
Here and below an occupied cube means a cube of positive mass for the
specified class.  In particular, for occupied $P\subset Q$ at depths
$b\ge a$,
\begin{equation}\label{reg:conditional-mass}
\frac{\rho_\alpha(P)}{\rho_\alpha(Q)}
 =R^{-[m_\alpha(b/N)-m_\alpha(a/N)]\pm C_n\lambda}.
\end{equation}
Constants depend only on $n$; the threshold for $N$ can depend on
$\lambda$.
\end{lemma}

\begin{proof}
Discard depth-$N$ cubes of mass less than $R^{-n-2}$.  There are at most
$R^n$ such cubes, so this costs at most $R^{-2}$.  Choose depths
\[
0=j_0<j_1<\cdots<j_J=N,
\qquad j_i-j_{i-1}\le\lambda N+1,
\qquad J\le C/\lambda.
\]
Give each remaining leaf its dyadic mass bin as an initial label.  There
are $O_n(N)$ possible initial labels.  Work upwards through the selected
depths.  If $\ell$ is a label at depth $j_i$, count, within each
depth-$j_{i-1}$ parent, the depth-$j_i$ children carrying $\ell$; extend
$\ell$ by the dyadic bin of this positive count.  The extension is given
to all leaves of label $\ell$ in that parent.  A parent can carry
several labels, but the sets of leaves having different labels are
disjoint.  At every extension there are at most $O_n(N)$ possible count
bins.  Consequently the number of final labels is at most
$N^{C_n/\lambda}$, after adjusting constants for sufficiently large $N$.

Fix a final label.  For each $i$ there is a positive integer $b_i$,
which is a power of two, such that every occupied depth-$j_{i-1}$ parent
has between $b_i$ and $2b_i$ occupied depth-$j_i$ children in this label.
To check that these comparisons survive the later extensions, fix a
partial label and a node at the current selected depth. All leaves with
that label in this node receive the same extension, determined by its
parent and the old label. At each subsequent, higher step the extension
is again constant on this entire group. Backward induction through the
selected depths therefore shows that a final label retains or removes
each earlier child group as a whole; it never splits that group.
Consequently the branching comparisons hold simultaneously. Also
\[
1\le b_i\le 2^{n(j_i-j_{i-1})}.
\]
All leaf masses are in one dyadic bin.  Counting leaves below a selected
node, and then normalizing the final class, therefore gives
\[
\rho_\alpha(Q)=2^{\pm C J}
                  \Bigl(\prod_{h\le i}b_h\Bigr)^{-1}
\quad\text{at depth }j_i.
\]
Define
\[
m_\alpha(j_i/N)=\frac1N\sum_{h\le i}\log_2b_h
\]
and interpolate linearly.  Its slopes lie in $[0,n]$ exactly.  At an
intermediate depth, an occupied cube is contained in its preceding
selected-depth parent and contains an occupied descendant at the next
selected depth.  Both the possible mass variation and the profile
variation across this gap are at most $R^{n\lambda+O_n(1/N)}$.
The remaining comparison error is $2^{O(J)}$, whose base-$R$
logarithm is $O(1/(\lambda N))$.  For fixed $\lambda$ and sufficiently
large $N$ these errors imply Equation~\eqref{reg:cell-mass}.

Discard final classes of original mass less than $R^{-\lambda}$.
Their total mass is bounded by the number of labels times
$R^{-\lambda}$, as asserted.  For any retained class, normalization
multiplies an original upper ball bound by at most $R^\lambda$.
Covering a depth-$j$ cube by a bounded number of balls of radius
$2^{-j}$ gives
\[
\rho_\alpha(Q)\le C_nKR^\lambda2^{-sj}.
\]
Comparison with the lower bound in Equation~\eqref{reg:cell-mass}
proves Equation~\eqref{reg:profile-lower-bound} at integer depths,
with a larger dimensional constant.  Linear interpolation gives it at
every $u$.  Dividing two instances of Equation~\eqref{reg:cell-mass}
proves Equation~\eqref{reg:conditional-mass}.
\end{proof}

For use in scale-compactness arguments, the proof also gives the
following precise interpretation of subpower regularity.  If
$\lambda=\lambda_N\to0$ sufficiently slowly that
$\lambda_N^2N/\log N\to\infty$, then the number of classes, all mass
comparison losses, and the normalization losses of retained classes
are $R^{o(1)}$, while the discarded mass tends to zero.  Their
profiles remain exactly $n$-Lipschitz.  This
observation alone does not assert a fixed power of discarded mass:
fixed positive deletion exponents are obtained by choosing a fixed
small regularization parameter after the relevant compactness
argument.

\begin{lemma}[Density pruning and multiplication]\label{reg:pruning-multiplication}
Let $\rho$ be a source probability and $\nu$ a finite pin measure,
both on $\R^n$,
and let $0\le a_0<\cdots<a_J\le N$ be a fixed finite collection of
dyadic depths.  Suppose pin--source distances on the sets in question
are between fixed positive constants.  Let $G'$ be a measurable set of
pairs.  For each $j<J$, suppose there is a measurable set
$E_j\supset G'$ and a number $f_j\ge0$ such that every
$(x,z)\in G'$, with $z\in Q$ at depth $a_j$, satisfies
\begin{equation}\label{reg:local-fibre}
\rho\{z'\in Q\cap(E_j)_x:
     |\pi_x(z')-\pi_x(z)|\le C2^{-a_{j+1}}\}
 \le f_j\rho(Q).
\end{equation}
The constant $C$ may be any sufficiently large fixed constant,
depending on the geometry and on the fixed angular width to be used
below.  For every $e>0$ one can delete pair mass at most
$(J+1)R^{-e}\nu(\R^n)$ from $G'$ and obtain $G''$ such that, for each
pin $x$, each initial cube $Q_0$ at depth $a_0$, and every spherical
ball $B$ of radius $C_12^{-a_J}$,
\begin{equation}\label{reg:multiplied-fibre}
\rho\{z\in Q_0\cap G''_x:\pi_x(z)\in B\}
 \le R^{eJ}\Bigl(\prod_{j=0}^{J-1}f_j\Bigr)\rho(Q_0).
\end{equation}
An identical conclusion holds without a lower pin--source separation
if Equation~\eqref{reg:local-fibre} instead tests the perpendicular
distance to the full line $x+\R(z-x)$ at width $C2^{-a_{j+1}}$,
provided pin and source sets remain bounded and pairs on the diagonal
have been removed.
\end{lemma}

\begin{proof}
For each endpoint depth delete the pairs $(x,z)\in G'$ whose source
cube $Q$ satisfies
\[
\rho(Q\cap G'_x)<R^{-e}\rho(Q).
\]
For a fixed pin, the total source mass deleted at one depth is at most
$R^{-e}$, by summing over the disjoint cubes.  Integration in the pin
and a union bound prove the deletion estimate.  Notice that the
density condition concerns $G'$, before this last deletion.  No
density assertion for $G''$ is needed.

Fix $x,Q_0,B$.  Call a cube at any endpoint depth hit if it contains a
point of $Q_0\cap G''_x$ whose direction lies in $B$.  In a hit parent
$Q$ at depth $a_j$, choose one such point $z_Q$.  If a child $P$ at
depth $a_{j+1}$ is hit, choose a corresponding point $z_P$.  The map
$\pi_x$ is uniformly Lipschitz on the separated source set, so every
$z'\in P\cap G'_x$ has
\[
|\pi_x(z')-\pi_x(z_Q)|
 \le C_22^{-a_{j+1}}+2C_12^{-a_J}
 \le C2^{-a_{j+1}}.
\]
As $G'\subset E_j$, Equation~\eqref{reg:local-fibre} and the
pre-deletion density of every hit child imply
\[
\sum_{\substack{P\subset Q\\P\text{ hit at }a_{j+1}}}\rho(P)
 \le R^e\sum_{P\text{ hit}}\rho(P\cap G'_x)
 \le R^ef_j\rho(Q).
\]
Iterate this inequality from $Q_0$ to depth $a_J$.  The mass on the
left of Equation~\eqref{reg:multiplied-fibre} is at most the sum of the
full masses of the final hit cubes, proving that Equation.

For the line-width variant, the two hit witnesses have directions
within $2C_12^{-a_J}$.  Bounded pin--source distances from above imply
that $z_P$ lies within $O(2^{-a_J})$ of the line through $x,z_Q$;
every point of $P$ is then within $O(2^{-a_{j+1}})$ of that line.
The preceding argument applies unchanged.  This step uses no
Lipschitz bound for the radial map near its pin.
\end{proof}

All selections above can be made measurably: the cube partitions and
label sets are finite, and the fibre integrals are measurable.
Enlarging fixed test widths permits the use of closed caps or closed
tubes without altering any exponent.
